% Tulkinta taulukoille, jotka tuotettiin komennolla Rscript paperi/laske_taulukot.R.
\paragraph{Tulkinta.}
Taulukoiden järjestys vastaa lauseita. Oletusasetuksilla
\pol{tuote} oli \pol{kulku}a parempi 193 siemenellä 200:sta, ja parittaisen
erotuksen luottamusväli $[0{,}38;\ 0{,}44]$ ei sisällä nollaa
(lause~\ref{la:aito}). Asetelmissa $G = 1$ ja $\eta = \varepsilon$ erotus oli
jokaisella siemenellä täsmälleen nolla, kuten lause~\ref{la:valttamaton}
ennustaa.

Suurin osa \pol{tuote}en edusta \pol{aikataulu}un nähden tulee kulkemisen
huomioimisesta: \pol{kulku} vähentää odotettuja kaatumisia noin 2,1 ja
tärinähavainnot vielä noin 0,4. Tärinähavainnot saavuttavat noin neljänneksen
täydellisen tiedon edusta. Oletusarvolla $\rho = 0$ liukkaus arvotaan
uudelleen keskiyöllä. Klo 00 kierroksella kuluvasta päivästä ei siksi ole
lainkaan tietoa, ja klo 06 kierroksella vain hiljaisen yön verran.
Informatiivisia ovat lähinnä klo 12 ja 18 kierrokset. Kun liukkaus säilyy
päivästä toiseen ($\rho = 0{,}9$), tieto kertyy, ja osuus nousee noin
60 prosenttiin.

Myös oraakkeli vähentää kaatumisia vain noin 24 prosenttia, koska huolto
kattaa kerrallaan 2 aluetta 20:stä. Asetelma $H = 12$ on oletuksen~\ref{ol:A}
ulkopuolella. Siinä \pol{tuote} oli edelleen keskimäärin parempi, mutta
lauseet eivät takaa tätä (huomautus~\ref{hu:H}).

Luvut ovat simuloituja kaatumisia $300 \times 300$ metrin alueella yhden kuukauden
aikana. Ne eivät ole vakuutusvahinkoja tai Turun todellisia tapahtumia.
